We'll start with taking a look at \textbf{Footprinting, Scanning} and 
\textbf{Enumeration}. \newline
\newline Attackers, before performing an attack, will perform 
some kind of investigation to determine the vulnerabilities
they can use at their own advantage. We are considering 
professional hackers, who don't want to be detected 
when they're entering the system. \newline 
Attackers behave as insiders when breaking a system, even 
if he could be outside the company. When we perform 
Ethical Hacking, we need to be inside the company, having
the same vision an hacker has. \newline 
Also, the major concern of an hacker is to be anonymous, 
as well as being stealty. At the basic of anonymity there 
is \texttt{Tor Project}. As the graphical user interface
for Tor, Vidalia is used, while for web filtering proxy
we have Privoxy. To solve IP addresses, it's possible to 
do this Tor as well, with \texttt{tor-resolve}. 
For proxychains, \texttt{Nmap} is well known. It's a tool 
that allows to do some sophisticated protocols. Once another 
tool is given by \texttt{nc} (netcat), used to send 
requests to servers. 

\subsubsection{Onion Router (TOR)}
We have the usual scheme, with Alice wanting to talk with 
Bob. However, they want the possibility to deny they're
communicating, leaving no trace of their communication.
The protocol works with three layers of routers, 
located around the world, and there are hundreds of thousand 
of them. Alice will pick three random routers. These 
routers will have special names: \texttt{guard node}, the 
node to which Alice is talking (entry point of the network);
\texttt{exit-node}. This is called Onion Router because 
the packet that is used is composed by many envelops. 
The main message is covered in many other messages. 
We have both encrypted and unencrypted links. 
Encryption is implemented by Tor protocol, but on top 
of this we can have more encryption on top of it. 
Tor protocol is also transparent to the communication. 
The packet being sent is encrypted three times, with three 
different keys, one after the other. The key is given 
by the Tor node that Alice has selected. Each node has 
a public and a private key. The guard node is the most 
external envelope of the encryption. The guard node then 
decrypts the outside of the message, forwarding it to the 
second node. \newline
So the sender prepares these onion messages, with every node 
removing one layer, and passing it to the next of the three 
nodes (randomly picked nodes). \newline
\textit{Correct this part, I got a bit lost}. 

\section{Footprinting}
So far, we've seen tools to find information about the 
target we attack. Before that, however, Footprinting is 
performed. It works without even sending to the target a 
single packet. So, we need to find information about the 
organization without even interacting. 
Footprinting should be executed also by us Ethical Hackers. 
This is internet Footprinting:
\begin{enumerate}
    \item Determine the scope of your activities: we need to 
    understand what I actually need to do, avoiding to 
    inadvertively commiting crimes. 
    \item Get proper authorization: the scope should be 
    linked with our authorization to the system. The contract 
    will tell us what we can and what we cannot do. 
    \item Check publicly available information
    \item WHOIS \& DNS enumeration 
    \item DNS interrogation 
    \item Network Reconnaissance.
\end{enumerate}


